% Generated draft: requires mathematical and source review.
\documentclass[11pt]{article}
\usepackage[margin=25mm]{geometry}
\usepackage{fontspec}
\setmainfont{Latin Modern Roman}
\usepackage{amsmath,amssymb,mathtools,unicode-math}
\setmathfont{Latin Modern Math}
\usepackage{enumitem,needspace,xurl,hyperref}
\setlength{\parindent}{0pt}
\setlength{\parskip}{5pt}
\setlength{\emergencystretch}{3em}
\newcommand{\N}{\mathbb{N}}
\newcommand{\Z}{\mathbb{Z}}
\newcommand{\Q}{\mathbb{Q}}
\newcommand{\R}{\mathbb{R}}
\newcommand{\C}{\mathbb{C}}
\newcommand{\F}{\mathbb{F}}
\newcommand{\A}{\mathbb{A}}
\newcommand{\PP}{\mathbb{P}}
\newcommand{\E}{\mathbb{E}}
\newcommand{\eps}{\varepsilon}
\newcommand{\dd}{\,\mathrm{d}}
\begin{document}
\section*{UnsolvedMath ID: 800554}
\textbf{Source problem code:} catalogue entry
\subsection{Definitions and mathematical statement}
For a graph \(G\), let \(\operatorname{cr}(G)\) denote its crossing number: the minimum number of edge crossings in a drawing of \(G\) in the plane, with crossings counted away from vertices under the usual general-position drawing conventions. The source also defines the genus of \(G\) as
\[
\gamma(G)=\min\{g\in\Z_{\ge 0}:G\text{ embeds in a surface of genus }g\}.
\]
Here an embedding has no edge crossings. The source does not specify whether surfaces are orientable or nonorientable; this convention must be fixed before \(\gamma\) is unambiguous.

Let \(\mathcal B_k\) denote the class called \emph{bar \(k\)-visibility graphs} in the supplied source. This is a placeholder for the source's graph class: the supplied record does not define a bar representation, the visibility and obstruction rules, or the permitted values of \(k\). Those missing conventions must be supplied before this class is precisely determined.

\textbf{Questions (Problem 28).} For each permitted value of \(k\), determine the largest crossing number and the largest genus of a graph in \(\mathcal B_k\). Once the class and surface convention have been fixed, the unrestricted extremal quantities can be expressed as
\[
C(k)=\sup\{\operatorname{cr}(G):G\in\mathcal B_k\},
\qquad
H(k)=\sup\{\gamma(G):G\in\mathcal B_k\},
\]
allowing \(+\infty\) if the corresponding values are unbounded. This supremum notation records the possibility that no finite largest value exists; it does not assert boundedness or attainment.

\textbf{Unresolved scope.} The supplied questions place no bound on the number of vertices and do not specify the graph domain, including finiteness and simplicity. If the intended question instead concerns graphs with a fixed number \(n\) of vertices, its extremal quantities would be
\[
C(n,k)=\sup\{\operatorname{cr}(G):G\in\mathcal B_k,\ |V(G)|=n\},
\qquad
H(n,k)=\sup\{\gamma(G):G\in\mathcal B_k,\ |V(G)|=n\}.
\]
The supplied record does not establish that this order-dependent formulation is intended. Consequently, the display above is only a possible clarification, not a replacement for the source questions.
\subsection{Short English statement}
Problem 28 asks how large the crossing number and genus of a bar \(k\)-visibility graph can be. The supplied registry labels the record open. A precise formulation still requires the definition of bar \(k\)-visibility, the range of \(k\), the graph and surface conventions, and clarification of whether the number of vertices is fixed.
\subsection{Sources}
\begin{itemize}
\item Local UnsolvedMath registry, permanent ID \(800554\), exact source problem code \texttt{catalogue entry}, title \emph{Visibility Graphs of Points, Segments and Polygons --- Problem 28}. The supplied record attributes the problem to Ghosh and Goswami, with proposed year 2010, and labels its status \texttt{open}; that label is supplied curation, not an independently verified mathematical or literature assessment.
\item Ghosh and Goswami, \emph{Unsolved Problems in Visibility Graphs of Points, Segments and Polygons} (2010), Problem 28. Supplied source URL: \url{https://arxiv.org/abs/1012.5187}. The background reports access on 2026-10-16 and extraction by \texttt{arxiv\_latex}; these provenance details are supplied and were not independently verified here.
\item The supplied statement includes the crossing-number question inside an \texttt{unsolved} environment and then gives the genus definition and genus question immediately afterward. Both questions are retained in this draft.
\end{itemize}
\end{document}
